\section*{第 \ref{ch:b2:funcseq} 章 --- 函数列与
函数项级数}

\begin{solution}{exo:b2:funcseq:1}
$f_n(x) = \frac{x}{1 + nx}$：逐点极限为 $0$（在 $\intcc{0}{1}$ 上）。
一致收敛：$f_n$ 在 $\intcc{0}{1}$ 上递增（导数
$\frac{1}{(1+nx)^2} > 0$），故 $\norm{f_n}_\infty = f_n(1) =
\frac{1}{1+n} \to 0$，即 \emph{一致收敛}（在 $\intcc{0}{1}$ 上）。

$g_n(x) = nx\,\eu^{-nx^2}$：逐点极限为 $0$（指数
衰减胜过多项式增长）。上确界：$g_n' = n\eu^{-nx^2}(1 - 2nx^2)$ 的零点为 $x_n =
\frac{1}{\sqrt{2n}}$，此时 $g_n(x_n) = \sqrt{\frac n2}\,\eu^{-1/2}
\to \infty$，故在 $\intcc{0}{1}$ 上不一致收敛 --- 但在
$\intco{\delta}{1}$ 上一致收敛，因为其上 $g_n(x) \leq n\,\eu^{-n\delta^2}
\to 0$。

$h_n(x) = x^n(1 - x^n)$：逐点极限为 $0$（在 $\intcc{0}{1}$ 上）
（两个因子皆然；在 $x = 1$ 处 $h_n = 0$）。上确界：令 $u = x^n \in
\intcc{0}{1}$，$u(1-u) \leq \frac14$ 在 $u = \frac12$ 处取到，
即 $x = 2^{-1/n} \in \intoo{0}{1}$，故 $\norm{h_n}_\infty =
\frac14 \not\to 0$，故在 $\intcc{0}{1}$ 上不一致收敛；在
$\intcc{0}{a}$ 上一致收敛（$\sup \leq a^n \to 0$）。
\end{solution}

\begin{solution}{exo:b2:funcseq:2}
$\int_0^1 nx\,\eu^{-nx^2}\dd x = \bigl[-\tfrac12
\eu^{-nx^2}\bigr]_0^1 = \frac{1 - \eu^{-n}}{2} \to \frac12$，而
$\int_0^1 \lim g_n = 0$。无矛盾：
\cref{thm:b2:funcseq:integration} 要求在该线段上 \emph{一致收敛}，
而此处不成立（高度为
$\sim\sqrt n$ 的隆起向 $0$ 滑动）。
\end{solution}

\begin{solution}{exo:b2:funcseq:3}
在 $\intcc{-1}{1}$ 上正规收敛：$\norm{x^n/n^2}_\infty =
\frac{1}{n^2}$，可和，因此 $S$ 在该处连续
（由\cref{thm:b2:funcseq:seriestransfer}可知）。

求导：导出级数 $\sum \frac{x^{n-1}}{n}$ 在
每个 $\intcc{-a}{a}$ 上正规收敛（$a < 1$；$\sup =
\frac{a^{n-1}}{n}$）：$S$ 为 $C^1$ 函数，定义在 $\intoo{-1}{1}$ 上，且
\[
S'(x) = \sum_{n\geq1} \frac{x^{n-1}}{n} = \frac1x \sum_{n\geq1}
\frac{x^n}{n} = -\frac{\ln(1 - x)}{x}
\qquad (0 < \abs x < 1),
\]
最后一个等式即本科一年级学过的对数级数（
在\cref{ch:b2:powerseries}中严谨地重新推导）。
\end{solution}

\begin{solution}{exo:b2:funcseq:4}
固定 $x > 0$，该级数为交错级数，其中 $\frac{1}{n + x}
\downarrow 0$，故逐点收敛，且余项估计
$\abs{R_N(x)} \leq \frac{1}{N + 1 + x} \leq \frac{1}{N+1}$ 这一估计
\emph{一致}成立于 $\intco{\delta}{\infty}$（实际上在
$\intoo{0}{\infty}$ 上亦然）：一致收敛。（不是正规收敛：
$\norm{\frac{(-1)^n}{n+x}}_\infty = \frac{1}{n + \delta}$，
发散）。连续性由
\cref{thm:b2:funcseq:seriestransfer}可得。

函数方程：对 $F(x + 1)$ 作指标替换 $m = n + 1$：
\[
F(x+1) = \sum_{n\geq0} \frac{(-1)^n}{n + 1 + x}
= \sum_{m\geq1}\frac{(-1)^{m-1}}{m+x} ,
\]
于是，把 $m = 0$ 的项从 $F(x)$ 中分离出，
\[
F(x) + F(x+1)
= \frac{1}{x} + \sum_{m\geq1}
\frac{(-1)^m + (-1)^{m-1}}{m+x} = \frac1x .
\]
\end{solution}

\begin{solution}{exo:b2:funcseq:5}
设 $g_n = f_n - f \geq 0$（由假设，关于 $n$ 递减；
极限逐点为 $0$）；每个 $g_n$ 连续。任取 $\varepsilon
> 0$，令 $U_n = \{x : g_n(x) < \varepsilon\}$，它是开集（
开集的原像），单调递增（$g_{n+1} \leq g_n$），且覆盖 $K$
（由逐点收敛）。由博雷尔--勒贝格
（\cref{thm:b2:metric:borellebesgue}），有限个 $U_{n_1}
\subseteq \dots \subseteq U_{n_k}$ 覆盖 $K$，于是 $K = U_{n_k}$，
即 $\norm{g_{n_k}}_\infty \leq \varepsilon$，由单调性
$\norm{g_n}_\infty \leq \varepsilon$ 对一切 $n \geq n_k$ 成立，故一致
收敛。（单调性至关重要：
\cref{exo:b2:funcseq:2}中的滑动凸起在紧集上
逐点收敛，但无一致性。）
\end{solution}

\begin{solution}{exo:b2:funcseq:6}
作替换 $u = nx$：
\[
\int_0^1 \frac{n f(x)}{1 + n^2x^2}\dd x
= \int_0^n \frac{f(u/n)}{1 + u^2}\,\dd u .
\]
被积函数 $h_n(u) = \frac{f(u/n)}{1+u^2}\mathbf{1}_{u \leq n}$
逐点收敛于 $\frac{f(0)}{1+u^2}$（由 $f$ 在
$0$ 处连续），且被 $\frac{\norm f_\infty}{1 + u^2}$ 控制，
在 $\intco{0}{\infty}$ 上可积：由控制收敛
（\cref{thm:b2:integration:dominated}）得极限为
\[
\int_0^\infty \frac{f(0)}{1 + u^2}\dd u = \frac{\pi}{2} f(0) .
\]
（核集中于 $0$ 处，构成一个逼近恒等式。）
\end{solution}

\begin{solution}{exo:b2:funcseq:7}
对 $f(x) = x^2$，利用证明中的
第二组二项恒等式：$\sum_k k^2 p_k = n(n-1)x^2 + nx$。于是
\[
B_n(f)(x) = \sum_k \frac{k^2}{n^2}\,p_k
= \frac{n(n-1)x^2 + nx}{n^2}
= x^2 + \frac{x(1 - x)}{n} :
\]
$B_n(f) - f = \frac{x(1-x)}{n}$，其上确界范数为 $\frac{1}{4n} =
O\bigl(\frac1n\bigr)$，与预期一致。
\end{solution}

\begin{solution}{exo:b2:funcseq:8}
取 $\varepsilon = 1$，则存在 $N$，使得 $\norm{P_n - P_m}_{\infty,
\R} \leq 1$ 对所有 $m, n \geq N$ 成立。在 $\R$ 上有界的多项式
必为常数（非常数者趋于 $\pm\infty$）：$P_n - P_m =
c_{n,m}$ 为常数。于是对 $n \geq N$，$P_n = P_N + c_n$，其中
$c_n = P_n(0) - P_N(0)$ 收敛（$0$ 处逐点收敛）。
因此 $f = \lim P_n = P_N + \lim c_n$，是一个多项式。
\end{solution}

\begin{solution}{exo:b2:funcseq:9}
(a) $\norm{(3/4)^n\varphi(4^n\cdot)}_\infty = \frac12
(3/4)^n$：正规收敛，故 $W$ 连续
（见\cref{thm:b2:funcseq:seriestransfer}）。

(b) 固定 $x$、$m$；选取 $h_m = \pm\frac12 4^{-m}$ 的符号，使
线段 $\intcc{4^mx}{4^m(x + h_m)}$（长为 $\frac12$）
不含半整数，从而 $\varphi$ 在其上为斜率 $\pm1$ 的仿射映射
（这是可以做到的：长为 $\frac12$ 的区间至多含一个
半整数点，取避开它的那一侧）。

对 $n > m$，$4^n h_m = \pm\frac12 4^{n-m}$ 为整数，且
$\varphi$ 以 $1$ 为周期：差中第 $n$ 项
为零。

对 $n = m$，$\abs{\varphi(4^m x + 4^m h_m) - \varphi(4^m x)} =
\abs{4^m h_m} = \frac12$（$\varphi$ 在该
线段上为斜率 $\pm 1$ 的仿射映射），故该项对差商的贡献恰为 $(3/4)^m \cdot
\frac{1/2}{\abs{h_m}} = (3/4)^m\,4^m = 3^m$。

对 $n < m$，$1$-利普希茨的 $\varphi$ 给出
$\bigl|(3/4)^n\bigl(\varphi(4^nx + 4^nh_m) -
\varphi(4^nx)\bigr)\bigr| \leq (3/4)^n 4^n\abs{h_m} =
3^n\abs{h_m}$：每项对差商的贡献至多为 $3^n$。

于是
\[
\Bigl|\frac{W(x + h_m) - W(x)}{h_m}\Bigr|
\geq 3^m - \sum_{n=0}^{m-1} 3^n
= 3^m - \frac{3^m - 1}{2} = \frac{3^m + 1}{2}
\longrightarrow \infty .
\]
若 $W$ 在 $x$ 处可导，则沿
$h_m \to 0$ 的每个差商都应收敛于 $W'(x)$，矛盾。$W$ 是
处处连续，处处不可导。
\end{solution}

\begin{solution}{exo:b2:funcseq:10}
逐点收敛：对 $x \in \intco{0}{1}$，该级数是公比
为 $-x$ 的几何级数，
\[
\sum_{n\geq0}(-1)^n x^n(1-x) = \frac{1-x}{1+x},
\]
而在 $x = 1$ 处每一项均为零，和为 $0 = \frac{1-1}{2}$，
与之相符 --- 和函数在 $\intcc{0}{1}$ 上连续。
一致性：余项为几何级数的尾项，
\[
\abs{R_N(x)} = \frac{x^{N+1}(1-x)}{1+x} \leq x^{N+1}(1 - x)
\leq \max_{\intcc01} t^{N+1}(1-t)
= \frac{1}{N+2}\Bigl(\frac{N+1}{N+2}\Bigr)^{\!N+1}
\leq \frac{1}{N+2} \to 0 ,
\]
对 $x$ 一致成立。不是正规收敛：$\norm{u_n}_\infty =
\max x^n(1-x) = \frac{1}{n+1}\bigl(\frac{n}{n+1}\bigr)^n \sim
\frac{1}{\eu\,n}$，且 $\sum \frac1{\eu n}$ 发散。对比：
$\sum x^n(1-x)$ 的部分和为 $1 - x^{N+1}$，
逐点收敛于\emph{不连续}的 $\mathbf 1_{\intco01}$：由
\cref{thm:b2:funcseq:continuity}，该收敛不可能
在 $\intcc{0}{1}$ 上一致。
\end{solution}

\begin{solution}{exo:b2:funcseq:11}
极限函数 $f$ 连续（见\cref{thm:b2:funcseq:continuity}）。
于是
\[
\abs{f_n(x_n) - f(x)}
\leq \abs{f_n(x_n) - f(x_n)} + \abs{f(x_n) - f(x)}
\leq \norm{f_n - f}_\infty + \abs{f(x_n) - f(x)} ,
\]
两项都趋于 $0$（由一致收敛性及
$f$ 在 $x$ 处的连续性）。若仅逐点收敛，则有反例：
$g_n(x) = nx\,\eu^{-nx^2} \to 0$ 在 $\intcc{0}{1}$ 上逐点收敛，
且 $g_n$ 与极限函数均连续，但在 $x_n =
\frac{1}{\sqrt{2n}} \to 0$ 处：
\[
g_n(x_n) = \sqrt{\frac n2}\,\eu^{-1/2} \longrightarrow +\infty
\neq 0 = f(0) .
\]
\end{solution}

\begin{solution}{exo:b2:funcseq:12}
\begin{enumerate}
  \item 归纳法。$n = 1$ 时即为定义。假设
        对 $n$ 的公式成立，令 $g(x) = \int_0^x \frac{(x-t)^n}{n!}
        f(t)\dd t$。被积函数在 $(x,t)$ 上连续，
        且为 $C^1$（关于 $x$）时，变上限含参积分
        的求导公式如下：
        \[
        g'(x) = \frac{(x-x)^n}{n!}f(x)
        + \int_0^x \frac{(x-t)^{n-1}}{(n-1)!}f(t)\dd t
        = T^nf(x)
        \]
        （将 $g(x+h) - g(x)$ 拆为
        窄条积分 $\int_x^{x+h}$，它为 $O(h\cdot\sup)$，被积函数在
        $t = x$ 处按 $h^n$ 的量级消失，另一部分为固定
        积分对 $x$ 的增量部分，借助中值不等式
        及连续性估计）。另外 $(T^{n+1}f)' = T^nf$
        （微积分基本定理），且 $g(0) =
        T^{n+1}f(0) = 0$：$T^nf$ 的两个原函数，
        在 $0$ 处为零者必相同，故 $T^{n+1}f = g$。其估计为：
        \[
        \abs{T^nf(x)} \leq \norm f_\infty
        \int_0^x \frac{(x-t)^{n-1}}{(n-1)!}\dd t
        = \norm f_\infty\,\frac{x^n}{n!}
        \leq \frac{\norm f_\infty}{n!} .
        \]
  \item $\sum_n \norm{T^nf}_\infty \leq \eu\,\norm f_\infty$：
        正规收敛，从而一致收敛；$S$ 连续。
        部分和满足 $S_N = f + T S_{N-1}$，且 $T$
        是 $1$-利普希茨的（范数为 $\norm\cdot_\infty$）
        （$\abs{Tg(x)} \leq x\norm g_\infty$）：令 $N \to
        \infty$ 于两端，得 $S = f + TS$。
  \item 令 $V(x) = f(x) + \eu^x\int_0^x \eu^{-t}f(t)\dd t$。
        于是 $V - f$ 为 $C^1$，且 $(V-f)'(x) =
        \eu^x\int_0^x\eu^{-t}f + f(x) = V(x)$，又 $(TV)' = V$
        且 $(V - f)(0) = TV(0) = 0$：故 $V - f = TV$，即
        $V$ 满足该方程。唯一性：若 $S_1, S_2$ 为
        连续解，$D = S_1 - S_2$ 满足 $D =
        TD$，故 $D = T^nD$ 对一切 $n$ 成立，且 $\norm D_\infty
        \leq \frac{\norm D_\infty}{n!} \to 0$：$D = 0$。
        因此
        \[
        \sum_{n\geq0} T^nf(x) = f(x) +
        \int_0^x \eu^{x-t}f(t)\,\dd t .
        \]
        （级数 $\sum T^n$ 是算子的几何级数，
        初步体现了预解式
        $(\mathrm{Id} - T)^{-1}$ 的思想，将在本科三年级
        分册中展开）。
\end{enumerate}
\end{solution}

\begin{solution}{pb:b2:funcseq:1}
\textbf{1.} 线性由公式易得。正性：
权重 $p_k(x) \geq 0$，故 $f \geq 0$ 蕴含 $B_nf \geq
0$；单调性只需对 $g - f$ 应用即可。界：$\pm f \leq
\norm f_\infty$ 可得 $\pm B_nf \leq \norm f_\infty B_ne_0 =
\norm f_\infty$。端点：$p_k(0) = \mathbf 1_{k=0}$ 与
$p_k(1) = \mathbf 1_{k=n}$，故 $B_nf(0) = f(0)$，$B_nf(1) =
f(1)$。

\textbf{2.} 将 $(x+y)^n = \sum_k\binom nk x^ky^{n-k}$ 关于 $x$ 求导，
两边乘以 $x$，并令 $y = 1 - x$：
\[
nx = \sum_k k\,p_k(x) ;
\]
再求一次导，并乘以 $x^2$：$n(n-1)x^2 = \sum_k k(k-1)p_k(x)$。
由此 $B_ne_0 = 1$（二项式定理），$B_ne_1(x) =
\frac{nx}{n} = x$，以及
\[
B_ne_2(x) = \frac{\sum_k k^2p_k}{n^2}
= \frac{n(n-1)x^2 + nx}{n^2}
= x^2 + \frac{x(1-x)}{n} .
\]

\textbf{3.} 展开：
\[
\sum_k\Bigl(\frac kn - x\Bigr)^{\!2} p_k
= B_ne_2(x) - 2x\,B_ne_1(x) + x^2
= \frac{x(1-x)}{n} .
\]
由柯西–施瓦茨不等式，取分解 $\abs{k/n - x}\sqrt{p_k}
\cdot \sqrt{p_k}$：
\[
\sum_k\Bigl|\frac kn - x\Bigr| p_k
\leq \Bigl(\sum_k\Bigl(\frac kn - x\Bigr)^2
p_k\Bigr)^{\!1/2}
= \sqrt{\frac{x(1-x)}{n}} \leq \frac{1}{2\sqrt n},
\]
利用 $x(1-x) \leq \frac14$。

\textbf{4.} 权重 $p_k(x)$ 非负且和为 $1$，
重心为 $\sum_k \frac kn p_k(x) = x$（第 2 问）。凸函数 $f$ 的有限
琴生不等式（由两点定义归纳，
见本科一年级分册）给出
\[
f(x) = f\Bigl(\sum_k \frac kn\,p_k\Bigr)
\leq \sum_k f\Bigl(\frac kn\Bigr)p_k = B_nf(x) .
\]

\textbf{5.} 在 $\{k : \abs{k/n - x} > \delta\}$ 上有
$\bigl(\frac{k/n - x}{\delta}\bigr)^2 > 1$，于是
\[
\sum_{\abs{k/n-x}>\delta} p_k
\leq \frac{1}{\delta^2}\sum_k\Bigl(\frac kn -
x\Bigr)^2p_k
= \frac{x(1-x)}{n\delta^2} \leq \frac{1}{4n\delta^2} .
\]
解读：$p_k(x)$ 是样本频率 $S_n/n$ 的分布，
即 $n$ 次偏置为 $x$ 的抛硬币试验；其均值为 $x$，方差为
$\frac{x(1-x)}n \to 0$，而上式即切比雪夫
不等式：质量集中于 $x$，故以之对 $f$ 作平均，
在极限下重现 $f(x)$
（\cref{ch:b2:genfun} 中将这些术语正式化）。

\textbf{6.} $\omega \leq 2M < \infty$；单调性显然
（上确界取于更大的集合）。海涅定理：$f$ 在紧集上连续，则
一致连续，这恰好表示 $\omega(\delta) \to 0$
当 $\delta \to 0^+$ 时。次可加性：若 $\abs{s - t} \leq
\delta_1 + \delta_2$，点 $u$ 在线段 $\intcc st$ 上，
且 $\min(\delta_1, \abs{s-t})$ 为其与 $s$ 的距离，满足
$\abs{s-u} \leq \delta_1$，$\abs{u-t} \leq \delta_2$，且
$\abs{f(s)-f(t)} \leq \abs{f(s)-f(u)} + \abs{f(u)-f(t)}$。
反复应用得 $\omega(p\delta) \leq p\,\omega(\delta)$，其中 $p \in
\N^*$；对 $\lambda > 0$，取 $p = \lceil\lambda\rceil \leq 1
+ \lambda$：$\omega(\lambda\delta) \leq \omega(p\delta) \leq
p\,\omega(\delta) \leq (1+\lambda)\omega(\delta)$.

\textbf{7.} 由 $\sum p_k = 1$ 得：
\[
\abs{B_nf(x) - f(x)}
= \Bigl|\sum_k\bigl(f(k/n) - f(x)\bigr)p_k\Bigr|
\leq \sum_k\omega\bigl(\abs{k/n - x}\bigr)p_k .
\]
对每个 $k$，取第 6 问中的 $\lambda = \abs{k/n - x}/\delta$
得 $\omega(\abs{k/n-x}) \leq \bigl(1 +
\frac{\abs{k/n-x}}\delta\bigr)\omega(\delta)$；对 $p_k$ 加权求和
即得主估计。

\textbf{8.} 代入第 3 问中的界：
\[
\abs{B_nf(x) - f(x)} \leq \Bigl(1 +
\frac{1}{2\delta\sqrt n}\Bigr)\omega(\delta),
\]
对 $x$ 一致成立；取 $\delta = n^{-1/2}$ 时括号为
$\frac32$：$\norm{B_nf - f}_\infty \leq
\frac32\omega(n^{-1/2}) \to 0$，由第 6 问（海涅定理）。此
即以速率重新证明了 \cref{thm:b2:funcseq:weierstrass}。

\textbf{9.} $L$ 利普希茨条件即 $\omega(\delta) \leq L\delta$：
速率为 $\frac{3L}{2\sqrt n}$。$\alpha$ 赫尔德条件即
$\omega(\delta) \leq C\delta^\alpha$：速率为 $\frac{3C}{2}
n^{-\alpha/2}$。

\textbf{10.} 利用 $k\binom{2m}k = 2m\binom{2m-1}{k-1}$：
\[
\sum_{k=m+1}^{2m}k\binom{2m}k
= 2m\sum_{j=m}^{2m-1}\binom{2m-1}{j}
= 2m\cdot 2^{2m-2},
\]
因为 $j \mapsto 2m-1-j$ 将 $\{m,\dots,2m-1\}$ 双射到
$\{0,\dots,m-1\}$，故和为 $2^{2m-1}$ 的一半。另有
$\sum_{k=m+1}^{2m}\binom{2m}k = \frac{2^{2m} -
\binom{2m}m}{2}$（同理由对称性）。于是
\[
\sum_{k=m+1}^{2m}(k-m)\binom{2m}k
= m\,2^{2m-1} - m\,\frac{2^{2m} - \binom{2m}m}{2}
= \frac m2\binom{2m}m .
\]
代换 $k \mapsto 2m-k$ 将满足 $k < m$ 的项
映到满足 $k > m$ 的项上（二项式系数相等，$\abs{k-m}$ 也相等）：
绝对值之和为单侧和的两倍，即 $m\binom{2m}m$。

\textbf{11.} 当 $x = \frac12$ 时，$p_k(\tfrac12) =
\binom{2m}k2^{-2m}$，且 $f(\tfrac12) = 0$：
\[
B_{2m}f\Bigl(\frac12\Bigr)
= \sum_k\Bigl|\frac{k}{2m} - \frac12\Bigr|
\binom{2m}k 2^{-2m}
= \frac{2^{-2m}}{2m}\,m\binom{2m}m
= \frac{\binom{2m}m}{2\cdot4^m}
\sim \frac{1}{2\sqrt{\pi m}}
\]
由\cref{ex:b2:comparison:centralbinomial}可得。由于 $\omega_f
(\delta) = \delta$ 在此成立（函数为 $1$-利普希茨的，且
界可取到），第 8 题预言至多为
$\frac32(2m)^{-1/2}$：真实误差 $\frac{1}{2\sqrt{\pi m}}$
恰好与 $n^{-1/2}$ 同阶 --- 速率在相差
一个常数因子的意义下是精确的。

\textbf{12.} 由 $f \leq g$ 得 $g - f \geq 0$，故 $L(g-f) \geq
0$，即 $Lf \leq Lg$。由 $-\abs f \leq f \leq \abs f$：
$-L\abs f \leq Lf \leq L\abs f$，即 $\abs{Lf} \leq
L\abs f$。

\textbf{13.} 由海涅定理，取 $\delta$，使得 $\abs{f(s)-f(x)}
\leq \varepsilon$，只要 $\abs{s-x} \leq \delta$。若
$\abs{s - x} > \delta$，则 $\frac{(s-x)^2}{\delta^2} > 1$，且
$\abs{f(s)-f(x)} \leq 2M \leq \frac{2M}{\delta^2}(s-x)^2$。两种
情形下所要证的估计均成立。

\textbf{14.} 固定 $x$，由第 13 题，视为 $s$ 的函数有：
\[
-\varepsilon e_0 - \frac{2M}{\delta^2}q_x
\;\leq\; f - f(x)e_0
\;\leq\; \varepsilon e_0 + \frac{2M}{\delta^2}q_x,
\qquad q_x = e_2 - 2x\,e_1 + x^2e_0 .
\]
对单调线性算子 $L_n$ 应用（第 12 题），并在
$x$ 处取值：
\[
\abs{L_nf(x) - f(x)L_ne_0(x)}
\leq \varepsilon L_ne_0(x)
+ \frac{2M}{\delta^2}\bigl(L_ne_2(x) - 2xL_ne_1(x)
+ x^2L_ne_0(x)\bigr) .
\]

\textbf{15.} 记 $\alpha_j = L_ne_j - e_j$，则
$\norm{\alpha_j}_\infty \to 0$。由于 $e_2(x) - 2xe_1(x) +
x^2e_0(x) = 0$，有：
\[
L_ne_2(x) - 2xL_ne_1(x) + x^2L_ne_0(x)
= \alpha_2(x) - 2x\,\alpha_1(x) + x^2\alpha_0(x),
\]
其上确界范数至多为 $\norm{\alpha_2} + 2\norm{\alpha_1} +
\norm{\alpha_0} \to 0$。又 $L_ne_0 \to e_0$ 一致收敛，故
$L_ne_0 \leq 2$（对充分大的 $n$），且 $\abs{f(x)}\abs{L_ne_0(x) -
1} \leq M\norm{\alpha_0} \to 0$。结合第 14 题，
当 $n$ 充分大时，关于 $x$ 一致地有：
\[
\abs{L_nf(x) - f(x)} \leq 2\varepsilon +
\frac{2M}{\delta^2}\,o(1) + M\,o(1) \leq 3\varepsilon :
\]
$L_nf \to f$ 一致收敛 --- 科罗夫金定理。

\textbf{16.} $B_ne_0 = e_0$ 与 $B_ne_1 = e_1$ 精确成立，且
$\norm{B_ne_2 - e_2}_\infty = \max_x\frac{x(1-x)}{n} =
\frac{1}{4n} \to 0$（第 2 题）：科罗夫金定理适用，且
由此第三次得到魏尔斯特拉斯定理。

\textbf{17.} $I_nf$ 关于 $f$ 是线性的（节点值如此），且在
每个小区间上，非负节点值的仿射插值为
非负，即正性。$I_ne_0 = e_0$ 与 $I_ne_1 = e_1$ 成立，
因仿射函数等于它自身的插值。在
小区间 $\intcc ab$ 上（$b - a = \frac1n$），
$e_2$ 的仿射插值为 $L(t) = (a+b)t - ab$，且
\[
L(t) - t^2 = (t-a)(b-t) \in
\intcc{0}{\tfrac{(b-a)^2}{4}} ,
\]
其最大值在中点处取得：$\norm{I_ne_2 - e_2}_\infty =
\frac{1}{4n^2} \to 0$。由科罗夫金定理，$I_nf \to f$ 一致地
对每个连续的 $f$ 成立 --- 折线逼近，无需
进一步的估计。

\textbf{18.} 给定 $f$ 与 $\varepsilon$：魏尔斯特拉斯定理给出
多项式 $P = \sum_{j=0}^d a_jx^j$，满足 $\norm{f -
P}_\infty \leq \frac\varepsilon2$；将每个 $a_j$ 替换为
有理数 $b_j$，使 $\abs{a_j - b_j} \leq
\frac{\varepsilon}{2(d+1)}$，则
$\intcc{0}{1}$ 上的上确界范数至多改变 $\frac\varepsilon2$。有理系数多项式的集合
是可数个可数集之并（对
$d$），因而可数，且稠密：
$C(\intcc{0}{1})$ 是可分的。

\textbf{19.} 由线性，$\int_0^1 fP = 0$ 对每个
多项式 $P$ 成立。选取多项式 $P_n \to f$ 一致收敛
（由魏尔斯特拉斯定理）：
\[
\Bigl|\int_0^1 f^2\Bigr|
= \Bigl|\int_0^1 f\,(f - P_n)\Bigr|
\leq \norm f_\infty\,\norm{f - P_n}_\infty
\longrightarrow 0 ,
\]
于是 $\int_0^1 f^2 = 0$。若 $f(x_0) \neq 0$，由连续性得
$f^2 \geq c > 0$ 在某子区间上成立，与
积分为零矛盾，故 $f = 0$。因此两个连续函数若
具有相同的矩 $\int f t^n$，则必相等。

\textbf{20.} 记 $p_{n,k}(x) = \binom nk x^k(1-x)^{n-k}$，并
约定 $p_{n-1,-1} = p_{n-1,n} = 0$，由乘积法则，
及 $k\binom nk = n\binom{n-1}{k-1}$，$(n-k)\binom nk =
n\binom{n-1}{k}$，可得：
\[
p_{n,k}'(x) = n\bigl(p_{n-1,k-1}(x) - p_{n-1,k}(x)\bigr) .
\]
以 $f(k/n)$ 为权求和，并在第一个
和中平移指标（阿贝尔分部求和）：
\[
(B_nf)'(x) = n\sum_{j=0}^{n-1}\Bigl(f\Bigl(\frac{j+1}n\Bigr)
- f\Bigl(\frac jn\Bigr)\Bigr)p_{n-1,j}(x) .
\]

\textbf{21.} 由中值定理，$f(\frac{j+1}n) -
f(\frac jn) = \frac1n f'(\xi_j)$，其中 $\xi_j \in
\intoo{j/n}{(j+1)/n}$，于是 $(B_nf)'(x) = \sum_j
f'(\xi_j)\,p_{n-1,j}(x)$。节点 $\frac{j}{n-1}$ 也属于
$\intcc{j/n}{(j+1)/n}$（两个不等式都化为 $j \leq
n-1$），因此 $\abs{\xi_j - \frac j{n-1}} \leq \frac1n$，且
\[
\bigl|(B_nf)'(x) - B_{n-1}(f')(x)\bigr|
\leq \sum_j\Bigl|f'(\xi_j) -
f'\Bigl(\frac{j}{n-1}\Bigr)\Bigr| p_{n-1,j}(x)
\leq \omega_{f'}\Bigl(\frac1n\Bigr) \longrightarrow 0
\]
一致地成立。由于 $B_{n-1}(f') \to f'$ 一致收敛
（由\cref{thm:b2:funcseq:weierstrass} 应用于连续函数
$f'$），由三角不等式得 $(B_nf)' \to f'$
一致收敛。此时多项式 $P_n = B_nf$ 收敛到 $f$
在 $C^1$ 意义下。

\textbf{22.} 在 $x$ 处的泰勒--拉格朗日公式：$f(\frac kn) - f(x) =
f'(x)(\frac kn - x) + \frac{f''(\xi_k)}2(\frac kn - x)^2$.
以 $p_k$ 为权求和，线性项消失（第 2 问）：
\[
\abs{B_nf(x) - f(x)}
\leq \frac{\norm{f''}_\infty}{2}\sum_k\Bigl(\frac kn -
x\Bigr)^2p_k
= \frac{\norm{f''}_\infty}{2}\cdot\frac{x(1-x)}{n}
\leq \frac{\norm{f''}_\infty}{8n} .
\]

\textbf{23.} 对 $(x+y)^n$ 再求两次导数，得到
阶乘矩，其中 $n_{(j)} = n(n-1)\cdots(n-j+1)$：
\[
\sum_k k_{(j)}\,p_k = n_{(j)}\,x^j \qquad (j = 3, 4),
\]
并利用 $k^3 = k_{(3)} + 3k_{(2)} + k$，$k^4 = k_{(4)} + 6k_{(3)} +
7k_{(2)} + k$ 将它们化为幂矩：
\[
\sum_k k^3p_k = n_{(3)}x^3 + 3n_{(2)}x^2 + nx,
\qquad
\sum_k k^4p_k = n_{(4)}x^4 + 6n_{(3)}x^3 + 7n_{(2)}x^2 + nx .
\]
展开 $(k - nx)^4$ 并合并同类项（这是一项耐心但纯属机械的计算，
借助四个幂矩）：
\[
\sum_k(k-nx)^4p_k = nx(1-x)\bigl(1 + 3(n-2)x(1-x)\bigr) .
\]
由 $x(1-x) \leq \frac14$，右端至多为
$\frac n4\bigl(1 + \frac{3n}4\bigr) = \frac{3n^2}{16} +
\frac n4 \leq n^2$ 对 $n \geq 1$ 成立。

\textbf{24.} 泰勒公式的佩亚诺形式在 $x$ 处定义了 $\eta(t) =
\frac{f(t) - f(x) - f'(x)(t-x) - \frac12f''(x)(t-x)^2}
{(t-x)^2}$，当 $t \neq x$ 时，$\eta(x) = 0$：由泰勒--拉格朗日公式
$\eta(t) = \frac12\bigl(f''(\xi) - f''(x)\bigr)$，其中 $\xi$
介于 $t$ 与 $x$ 之间，于是 $\abs\eta \leq \norm{f''}_\infty$，且
$\eta(t) \to 0$，当 $t \to x$ 时（由 $f''$ 的连续性）。以
$p_k$ 为权对展开式求和，并利用第 2、3 问：
\[
n\bigl(B_nf(x) - f(x)\bigr)
= \frac{x(1-x)}{2}f''(x)
+ n\sum_k\eta\Bigl(\frac kn\Bigr)\Bigl(\frac kn -
x\Bigr)^2p_k .
\]
给定 $\varepsilon$，取 $\delta$ 使得 $\abs\eta \leq
\varepsilon$ 在 $\abs{t - x}\leq\delta$ 上成立。近部分至多为
$\varepsilon\,n\cdot\frac{x(1-x)}n \leq \varepsilon$。远部分：
取 $C = \norm{f''}_\infty$，并利用第 23 问，
\[
n\,C\sum_{\abs{k/n-x}>\delta}\Bigl(\frac kn -
x\Bigr)^2p_k
\leq \frac{nC}{\delta^2}\sum_k\Bigl(\frac kn -
x\Bigr)^4p_k
= \frac{nC}{\delta^2 n^4}\sum_k(k-nx)^4p_k
\leq \frac{C}{\delta^2 n} \longrightarrow 0 .
\]
由此 $n(B_nf(x) - f(x)) \to \frac{x(1-x)}2f''(x)$ ---
沃罗诺夫斯卡娅定理。对 $f = e_2$，此式对每个
$n$ 都精确成立（见\cref{exo:b2:funcseq:7}）：$\frac1n$ 这一精度上限是真实的。

\textbf{25.} (i) 正性把逐点不等式化为
算子不等式：它给出了范数界、琴生不等式、
切比雪夫不等式，以及科罗夫金定理的全部内容 --- 仅靠线性性在此
什么也证明不了。(ii) 紧性经由海涅定理进入（第
6、13 问），经由 $f$ 的有界性，以及范数
$\norm\cdot_\infty$ 本身有限这一事实。(iii) 三个检验函数
就足够了，因为正性把一切归结为控制
$L_n$ 在单一族 $(s-x)^2 = e_2 - 2xe_1 + x^2e_0$ 上的取值，
其张成空间与 $e_0, e_1, e_2$ 所张成的相同。(iv) 伯恩斯坦多项式
以真实的速率 $\omega(n^{-1/2})$ 对每个连续的 $f$ 收敛
（这是最佳的，见第 11 问），但饱和于 $\frac1n$，对光滑的 $f$
（见第 24 问）；傅里叶一章沿用同一纲领，
以费耶尔核处理周期函数 --- 另一个正算子，
具有同样的优点与同样的克制。
\end{solution}
